\section{Three-Aim Structure \& Technical Roadmap}
The core of the Chalk Talk is the three-aim structure. The table below outlines how to pitch these aims to satisfy both the fundamental biologists and the translational engineers in the search committee, using principles established in multicellular programming \cite{daniels2024programming}.

\subsection{Aim Breakdown and Feasibility Analysis}

\begin{table}[h]
\centering
\caption{Strategic Breakdown of Research Aims}
\label{tab:aims}
\begin{tabular}{@{}p{3.5cm}p{4.8cm}p{4.8cm}p{3.5cm}@{}}
\toprule
\textbf{Aim Title} & \textbf{Biological Hypothesis} & \textbf{Key Methodology} & \textbf{Funding \& Risk} \\ \midrule
\textbf{Aim 1: Self-Patterning Morphogen Gradients} & Localized synthetic morphogens can establish robust spatial patterns in stem cell layers. & Notch-based synNotch receptors coupled to fluorescent reporter readouts. & \textbf{Low Risk} \\ Nexa Career Grant \\
\textbf{Aim 2: Closed-Loop ECM Remodeling} & Matrix-degrading enzymes controlled by synthetic promoters prevent fibrous scarring. & Doxycycline-inducible synthetic promoters driving matrix metalloproteinases. & \textbf{Medium Risk} \\ Apex Health (R01) \\
\textbf{Aim 3: Programmable Vascularization} & Engineered cell-autonomous migration can form self-assembling vascular networks. & Optogenetic control of chemokine receptors (CXCR4 mutants). & \textbf{High Risk} \\ Vertex Consortium \\ \bottomrule
\end{tabular}
\end{table}

\subsection{Deep-Dive \& Board Representation of Aims}

\subsubsection{Aim 1: Self-Patterning Morphogen Gradients (Years 1--2)}
\begin{itemize}
    \item \textbf{The Pitch:} Focus on preliminary data from postdoc work. Show that we can control signaling range by tweaking ligand shed rate.
    \item \textbf{Board Details:} Sketch a 1D cell array. Use colored chalk to show morphogen emission from a source cell, diffusion across neighbors, and the resulting step-function activation thresholds.
\end{itemize}

\subsubsection{Aim 2: Closed-Loop ECM Remodeling (Years 2--4)}
\begin{itemize}
    \item \textbf{The Pitch:} This is the core R01-equivalent project. We will engineer cells to sense mechanical strain and secrete remodelers in response, preventing the fibrotic pathway typical of chronic wounds.
    \item \textbf{Board Details:} Sketch the closed-loop feedback block diagram. Input: mechanical stress; Controller: engineered promoter; Output: matrix metalloproteinases (MMPs) secretion.
\end{itemize}

\subsubsection{Aim 3: Programmable Vascularization (Years 3--5)}
\begin{itemize}
    \item \textbf{The Pitch:} The high-risk/high-reward component. Traditional scaffolds suffer from inner necrosis due to lack of vascular networks. By programming cell-autonomous vascular recruitment, we bypass the need for external printing.
    \item \textbf{Board Details:} Sketch a vascular lumen crossing a synthetic boundary. Write down the genetic circuit wiring that controls migration velocity.
\end{itemize}
